\begin{table*}[!htbp]
\centering
\scriptsize
\caption{Branch-conditional coverage (L2)}
\label{tab:4}
\vspace{0.5em}
\par\textbf{Panel A. Wide $\tau$=2, post-decision branches (main seed, 128 test events)}
\vspace{0.1em}

\begin{adjustbox}{max width=\linewidth}
\begin{tabular}{@{}lllllllllll@{}}
\toprule
Score & hold n & hold cov. & hold union & rebal n & rebal cov. & rebal union & rebal gain (SE) & overall width ratio & trivial bound & taxonomy hold / rebal \\
\midrule
raw & 104 & 0.9712 & 1.0000 & 24 & 0.6667 & 0.9167 & +0.250 (0.088) & 3.96$\times$ & 0.542 & cost\_only / clear\_gain \\
studentized & 104 & 0.9712 & 1.0000 & 24 & 0.6667 & 0.9583 & +0.292 (0.093) & 3.97$\times$ & 0.542 & cost\_only / clear\_gain \\
cqr & 104 & 0.9423 & 1.0000 & 24 & 0.7500 & 0.9167 & +0.167 (0.076) & 3.37$\times$ & 0.500 & clear\_gain / clear\_gain \\
\bottomrule
\end{tabular}
\end{adjustbox}
\vspace{0.5em}
\par\noindent\begin{minipage}{\linewidth}\raggedright{\scriptsize Notes. Nominal 0.90 ($\alpha=0.1$). The paired gain's standard error (one SE, not a 95\% interval) is in parentheses; a single-cell coverage SE is about 0.05--0.10 for the 24-event branch and 0.02--0.05 for the 104-event branch. \texttt{Trivial bound} is $1-(1-\hat c)/\hat q_b$ with $\hat c$ the measured overall marginal coverage and $\hat q_b=24/128$; it uses the observed $\hat c$, not the nominal level, so it is an empirical identity, not a guarantee. The branch-union threshold is $\max(\hat q_{\rm hold},\hat q_{\rm rebalance})$ on the calibration side; \texttt{post\_rebalance} is a post-decision label and is never used at prediction time. \texttt{width ratio} is overall union width / overall marginal width.}\end{minipage}\par

\vspace{0.4em}
\par\textbf{Panel B. Branch-union aggregate across 56 method$\times$direction blocks}
\vspace{0.1em}

\begin{adjustbox}{max width=\linewidth}
\begin{tabular}{@{}ll@{}}
\toprule
Statistic & Value \\
\midrule
Coverage dominance (union $\ge$ marginal) & 56/56 blocks, 112/112 branch cells \\
Nonempty calibration branches & 112/112 branch cells; smallest $n_b$ = 16 \\
Infinite branch thresholds ($\hat q_b=\infty$) & 0 cells \\
Taxonomy & cost\_only 44, clear\_gain 25, gain\_but\_fails 10, fails\_no\_clear\_gain 4, ambiguous 29 \\
Wide 81$\to$2 rebalance, identified studentized & 0.667 $\to$ 0.750 (width 1.27$\times$, ambiguous) \\
Wide 81$\to$2 rebalance, identified CQR & 0.750 $\to$ 0.875 (width 1.67$\times$, ambiguous) \\
\bottomrule
\end{tabular}
\end{adjustbox}
\vspace{0.5em}
\par\noindent\begin{minipage}{\linewidth}\raggedright{\scriptsize Notes. Taxonomy uses normal-approximation 95\% intervals per branch; the cells are descriptive, not simultaneous. \texttt{cost\_only} = marginal coverage already clearly at or above nominal; \texttt{clear\_gain} = a gain whose 95\% interval excludes zero, with union still possibly below nominal; \texttt{gain\_but\_fails} = a gain whose 95\% interval excludes zero, but union still clearly below nominal; \texttt{fails\_no\_clear\_gain} = below nominal with no gain whose 95\% interval excludes zero; \texttt{ambiguous} otherwise.}\end{minipage}\par

\vspace{0.4em}
\par\textbf{Panel C. Pre-decision Mondrian cannot isolate the post branch (text summary)}
\vspace{0.1em}

\par\noindent\begin{minipage}{\linewidth}\raggedright{\scriptsize Wide $\tau$=2, identified raw/studentized, three seeds, \texttt{onset\_action}/\texttt{onset\_side}: post-rebalance coverage ranges 0.583--0.708 across 12 cells, versus 0.971 on hold and 0.667 on rebalance under the unstratified marginal pipeline. Pre-decision stratification cannot recover the post-decision branch because the branch label is not in the pre-decision $\sigma$-field.}\end{minipage}\par

\vspace{0.4em}
\par\textbf{Panel D. Branch-level Wilson intervals (Wide $\tau$=2)}
\vspace{0.1em}

\begin{adjustbox}{max width=\linewidth}
\begin{tabular}{@{}lllll@{}}
\toprule
Score & Branch & $n$ & Marginal [95\% CI] & Union [95\% CI] \\
\midrule
raw & hold & 104 & 0.971 [0.919, 0.990] & 1.000 [0.964, 1.000] \\
raw & rebalance & 24 & 0.667 [0.467, 0.820] & 0.917 [0.742, 0.977] \\
studentized & hold & 104 & 0.971 [0.919, 0.990] & 1.000 [0.964, 1.000] \\
studentized & rebalance & 24 & 0.667 [0.467, 0.820] & 0.958 [0.798, 0.993] \\
cqr & hold & 104 & 0.942 [0.880, 0.973] & 1.000 [0.964, 1.000] \\
cqr & rebalance & 24 & 0.750 [0.551, 0.880] & 0.917 [0.742, 0.977] \\
\bottomrule
\end{tabular}
\end{adjustbox}
\vspace{0.5em}
\par\noindent\begin{minipage}{\linewidth}\raggedright{\scriptsize Notes. Wilson score 95\% intervals for branch coverage of the fitted score, conditional on the fitted quantile; they do not include calibration uncertainty. One branch cell at a time; they are not simultaneous. The rebalance cells have 24 test events, so their marginal intervals are about 0.35 wide. All three rebalance marginal intervals lie entirely below 0.90: the deficit is directional even at this sample size, but its magnitude is not pinned down. Every rebalance union interval contains 0.90: union coverage is compatible with nominal but does not establish nominal coverage. Panel B's aggregate taxonomy retains normal-approximation intervals for comparability across all 56 blocks; Panel D uses Wilson intervals because the normal approximation is unreliable at $n=24$.}\end{minipage}\par

\end{table*}
